% ============================================================================
%  CAPÍTULO III: DESARROLLO DEL TRABAJO
% ============================================================================

\section{Desarrollo del trabajo}

\subsection{Figura de ejemplo}

Figura con título y línea FUENTE al estilo APA (usa el logo institucional
incluido en \texttt{figuras/logo.png}):

\figuraAPA{fig:logo}{Logo institucional UPDS}{\includegraphics[width=0.35\textwidth]{logo.png}}{Universidad Privada Domingo Savio.}

\subsection{Captura pendiente}

Esta macro muestra un recuadro mientras la captura no exista en
\texttt{figuras/}; cuando guardes la imagen con ese nombre, aparece
automáticamente:

\capturaAPA{fig:captura-ejemplo}{Ejemplo de captura de pantalla}{captura-ejemplo.png}{Pantalla de ejemplo del sistema.}

\subsection{Bloque de código}

\begin{lstlisting}[language=bash]
# Comandos de ejemplo
sudo apt update
sudo systemctl restart nginx
\end{lstlisting}

\subsection{Terminal}

\begin{lstlisting}[style=terminal]
$ systemctl status nginx
nginx.service - A high performance web server
   Active: active (running)
\end{lstlisting}

\subsection{Diagrama de arquitectura de red}

Ejemplo de topología cliente–router–servidor (como en los informes reales):

\figuraAPA{fig:red}{Arquitectura de red del laboratorio}{%
  \begin{tikzpicture}[
      nodo/.style={draw, rounded corners, minimum width=2.4cm,
                   minimum height=0.9cm, align=center, font=\small},
      nube/.style={draw, ellipse, align=center, font=\small, minimum width=2.2cm},
      linea/.style={-Stealth, thick}]
    \node[nodo] (pc) {Cliente};
    \node[nodo, right=1.8cm of pc] (rtr) {Router\textbackslash 192.168.0.1};
    \node[nube, right=1.8cm of rtr] (net) {Internet};
    \node[nodo, below=1.4cm of rtr] (srv) {Servidor Debian};
    \draw[linea] (pc) -- (rtr);
    \draw[linea] (rtr) -- (net);
    \draw[linea] (srv) -- (rtr);
  \end{tikzpicture}%
}{Elaboración propia.}

\subsection{Diagrama de estados (autómata)}

Ejemplo de máquina de estados finitos con transiciones etiquetadas:

\figuraAPA{fig:fsm}{Diagrama de estados de ejemplo}{%
  \begin{tikzpicture}[
      estado/.style={draw, circle, minimum size=1cm, font=\small},
      linea/.style={-Stealth, thick},
      nodo/.style={align=center}]
    \node[estado] (q0) {$q_0$};
    \node[estado, right=1.8cm of q0] (q1) {$q_1$};
    \node[estado, accepting, right=1.8cm of q1] (q2) {$q_2$};
    \draw[linea] (q0) -- node[above] {0} (q1);
    \draw[linea] (q1) -- node[above] {1} (q2);
    \draw[linea] (q1) to[bend left] node[above] {0} (q0);
    % Si se rompe la línea siguiente, el estado inicial desaparece (nodos y
    % posicionamiento relativos: siempre definir los nodos antes de usarlos).
  \end{tikzpicture}%
}{Elaboración propia.}

\subsection{Diagrama de flujo}

\figuraAPA{fig:flujo}{Flujo de decisión de ejemplo}{%
  \begin{tikzpicture}[
      paso/.style={draw, rounded corners, minimum width=2.6cm, minimum height=0.9cm, align=center, font=\small},
      decision/.style={draw, diamond, aspect=2, align=center, font=\small, inner sep=1pt},
      linea/.style={-Stealth, thick}]
    \node[paso] (ini) {Inicio};
    \node[decision, below=1cm of ini] (d) {¿Servicio activo?};
    \node[paso, below=1cm of d] (ok) {Continuar};
    \node[paso, right=1.6cm of d] (err) {Reiniciar};
    \draw[linea] (ini) -- (d);
    \draw[linea] (d) -- node[right, font=\scriptsize] {sí} (ok);
    \draw[linea] (d) -- node[above, font=\scriptsize] {no} (err);
  \end{tikzpicture}%
}{Elaboración propia.}

\subsection{Casos que "siempre daban error" (prueba del fix)}

El bloque siguiente usa flechas con la sintaxis moderna
(\texttt{->}, \texttt{>=stealth}) y texto entre nodos. **Antes** del fix
de babel (\texttt{\\shorthandoff\{<>\}} + librería \texttt{babel}) esto
fallaba con \texttt{Argument of \\language@active@arg> has an extra \}}.
Si compila, el fix funciona:

\figuraAPA{fig:fix}{Prueba del fix babel/TikZ}{%
  \begin{tikzpicture}[nodo/.style={draw, rounded corners, minimum width=2cm,
                     minimum height=0.8cm, align=center}]
    \node[nodo] (a) {Antes};
    \node[nodo, right=2cm of a] (b) {Después};
    \draw[->, >=stealth, thick] (a) -- node[above, font=\scriptsize]{tcp/80} (b);
  \end{tikzpicture}%
}{Elaboración propia.}

{\footnotesize\textbf{Errores que ya no deben aparecer:} tokens activos de
babel-spanish (\texttt{<}, \texttt{>}) rompiendo estilos de flecha;
texto de nodos con acentos dentro de argumentos de macro (usar
\texttt{\\shorthandoff} ya activo); depender de librerías no cargadas
(automata: cargar \texttt{automata} si se usa \texttt{accepting};
arrows.meta ya está cargado).\par}
